% =============================================================================
%  Erasmo de Roterdam — De Stultitiae Laude
%  Aforismo medieval portado al sistema visual claudeeditorial.
%  Compilacion: xelatex aphorisms/erasmus-elogio-locura.tex
% =============================================================================
\documentclass{claudeeditorial}

\renewcommand{\DocKicker}{Aforismo · Anno Domini MDIX}
\renewcommand{\DocTitle}{Erasmvs Roterodamensis}
\renewcommand{\DocSubtitle}{De Stultitiae Laude · Sobre la arrogancia de los faltos de talento}
\renewcommand{\DocAuthor}{Patricio Hernández Ballester}
\renewcommand{\DocRole}{Aforismo}
\renewcommand{\DocDate}{Anno Domini MDIX}
\renewcommand{\DocRef}{APH-ERASMUS-001}
\renewcommand{\DocFooter}{claudeeditorial · aforismo}

\hypersetup{
  pdftitle={Erasmus - De Stultitiae Laude},
  pdfauthor={Patricio Hernandez Ballester}
}

\ClaudeRunningHeader

\begin{document}
\BodyCopy
\ClaudeCover

\ClaudeChapterPage{I}{De la vanidad de los ignorantes}{En el año del Señor de mil quinientos y nueve, el sabio Desiderio Erasmo de Róterdam escribió su célebre obra \emph{Stultitiae Laus}.}

\begin{claudehero}
\ClaudeLead{El humanista de Róterdam, con pluma sagaz y mordaz ironía, desenmascaró los vicios de su tiempo. \emph{El Elogio de la Locura} trasciende los siglos: en todo tiempo y lugar se hallan aquellos que, carentes de verdadero mérito, buscan aparentarlo mediante la soberbia y la presunción.}
\end{claudehero}

\section{La sentencia que trasciende los siglos}

Erasmo nos enseña una verdad cuyo eco aún resuena: \Highlight{«Cuanto menos talento tienen, más arrogancia, más vanidad, más orgullo»}.

Aunque estas palabras exactas no aparecen en forma literal en sus escritos, el espíritu de esta sentencia permea toda su crítica social, especialmente dirigida contra aquellos que:

\begin{itemize}
\item Se adornan con títulos pomposos para ocultar su ignorancia.
\item Recurren a la autoridad del cargo cuando les falta sabiduría.
\item Adoptan gestos teatrales y vestiduras rimbombantes.
\item Presumen de lo que no poseen ni comprenden.
\end{itemize}

\section{Las palabras del maestro}

\begin{claudequote}
Donde no hay virtud ni saber, se cubren las faltas con vestiduras rimbombantes, títulos pomposos y gestos teatrales.
\end{claudequote}

\begin{claudequote}
Los que no tienen talento suelen apoyarse en la autoridad del cargo o en la gravedad fingida.
\end{claudequote}

\section{La enseñanza perpetua}

Esta sabiduría erasmiana cobra nueva vida en cada época. Como reza el aforismo atribuido al gran humanista:

\begin{ClaudeTip}[Aforismo final]
Allí donde falta el talento, abunda la presunción.
\end{ClaudeTip}

\vfill

\begin{center}
{\UIFont\small\color{claudeStone}\itshape Scripta manent, verba volant\\
\smallskip
Lo escrito permanece, las palabras vuelan}
\end{center}

\ClaudeSignature{\DocAuthor}{\DocRole}

\end{document}
